
\documentclass{article}
\usepackage{colortbl}
\usepackage{makecell}
\usepackage{multirow}
\usepackage{supertabular}

\begin{document}

\newcounter{utterance}

\centering \large Interaction Transcript for game `eqbench', experiment `eq\_bench', episode 8 with reflexion{-}Qwen3.5{-}27B.
\vspace{24pt}

{ \footnotesize  \setcounter{utterance}{1}
\setlength{\tabcolsep}{0pt}
\begin{supertabular}{c@{$\;$}|p{.15\linewidth}@{}p{.15\linewidth}p{.15\linewidth}p{.15\linewidth}p{.15\linewidth}p{.15\linewidth}}
    \# & \multicolumn{2}{c}{Player} && \multicolumn{2}{c}{Game Master} \\
    \hline

    \theutterance \stepcounter{utterance}  
    & & & \multicolumn{4}{p{0.6\linewidth}}{
        \cellcolor[rgb]{0.9,0.9,0.9}{
            \makecell[{{p{\linewidth}}}]{
                \texttt{\tiny{[P1$\langle$GM]}}
                \texttt{You have attempted this same task before and did not succeed. These are the plans you wrote after those attempts.} \\
\texttt{Trial \#1: New plan: My previous attempt failed because I misidentified the primary emotional response to Elise's specific attack. When Elise countered Mike's concern by attacking his own character ("hides behind his books"), Mike's reaction is not primarily embarrassment (which implies shame) or anger (which implies aggression), but rather a defensive reaction to having his own vulnerabilities exposed. I incorrectly prioritized "Defensiveness" as the only high score, but the prompt asks for a relative intensity of *all* emotions. A more accurate strategy is to recognize that while Defensiveness is the dominant reaction (high score), the personal nature of the attack also triggers a significant amount of Anger (at the unfairness) and Embarrassment (at being exposed), while Indifference remains zero. Next time, I will analyze the specific rhetorical move (a counter{-}attack on character) and assign high scores to both Defensiveness and Anger, with a moderate score for Embarrassment, rather than isolating a single emotion. I will ensure the scores reflect the nuance that a character pushed to the brink of a personal argument feels a mix of these emotions, not just one.} \\
\\ 
\\ 
\\ 
\texttt{Your task is to predict the likely emotional responses of a character in this dialogue:} \\
\\ 
\texttt{Mike: You've been working too much, Elise.} \\
\texttt{Elise: And you've been reading too much. We all have our vices, Mike.} \\
\texttt{Mike: It's not a vice, it's a concern. You're burning yourself out.} \\
\texttt{Elise: Oh, so now you're my therapist? } \\
\texttt{Mike: No, I'm your friend. And as your friend, I'm worried about you.} \\
\texttt{Elise: Well, 'friend', I didn't ask for your worry.} \\
\texttt{Mike: You never do, Elise. That's the problem.} \\
\texttt{Elise: And what's that supposed to mean?} \\
\texttt{Mike: Nothing. Forget it.} \\
\texttt{Elise: No, say it. You've got something on your mind, spill it.} \\
\texttt{Mike: It's just... you don't let people in, Elise. You push everyone away.} \\
\texttt{Elise: And you're the expert on letting people in? You, who hides behind his books and his quiet little bookstore?} \\
\texttt{{[}End dialogue{]}} \\
\\ 
\texttt{At the end of this dialogue, Mike would feel...} \\
\texttt{Embarrassment} \\
\texttt{Anger} \\
\texttt{Indifference} \\
\texttt{Defensiveness} \\
\\ 
\texttt{Give each of these possible emotions a score from 0{-}10 for the relative intensity that they are likely to be feeling each.} \\
\\ 
\texttt{You must output in the following format, including headings (of course, you should give your own scores), with no additional commentary:} \\
\\ 
\texttt{Embarrassment: <score>} \\
\texttt{Anger: <score>} \\
\texttt{Indifference: <score>} \\
\texttt{Defensiveness: <score>} \\
\\ 
\\ 
\texttt{{[}End of answer{]}} \\
\\ 
\texttt{Remember: zero is a valid score, meaning they are likely not feeling that emotion. You must score at least one emotion > 0.} \\
\\ 
\texttt{Your answer:} \\
            }
        }
    }
     \\ \\

    \theutterance \stepcounter{utterance}  
    & \multicolumn{4}{p{0.6\linewidth}}{
        \cellcolor[rgb]{0.561,0.737,0.561}{
            \makecell[{{p{\linewidth}}}]{
                \texttt{\tiny{[P1$\rangle$GM]}}
                \texttt{Embarrassment: 6} \\
\texttt{Anger: 7} \\
\texttt{Indifference: 0} \\
\texttt{Defensiveness: 9} \\
            }
        }
    }
    & & \\ \\

    \theutterance \stepcounter{utterance}  
    & & & \multicolumn{2}{p{0.3\linewidth}}{
        \cellcolor[rgb]{0.9,0.9,0.9}{
            \makecell[{{p{\linewidth}}}]{
                \texttt{\tiny{[GM$|$GM]}}
                \texttt{Embarrassment: 6} \\
\texttt{Anger: 7} \\
\texttt{Indifference: 0} \\
\texttt{Defensiveness: 9} \\
            }
        }
    }
    & & \\ \\

    \theutterance \stepcounter{utterance}  
    & & & \multicolumn{2}{p{0.3\linewidth}}{
        \cellcolor[rgb]{0.9,0.9,0.9}{
            \makecell[{{p{\linewidth}}}]{
                \texttt{\tiny{[GM$|$GM]}}
                \texttt{Embarrassment: 4} \\
\texttt{Anger: 3} \\
\texttt{Indifference: 0} \\
\texttt{Defensiveness: 7} \\
            }
        }
    }
    & & \\ \\

    \theutterance \stepcounter{utterance}  
    & & & \multicolumn{2}{p{0.3\linewidth}}{
        \cellcolor[rgb]{0.9,0.9,0.9}{
            \makecell[{{p{\linewidth}}}]{
                \texttt{\tiny{[GM$|$GM]}}
                \texttt{game\_result = LOSE} \\
            }
        }
    }
    & & \\ \\

    \theutterance \stepcounter{utterance}  
    & & \multicolumn{4}{p{0.6\linewidth}}{
        \cellcolor[rgb]{0.95,0.95,0.95}{
            \makecell[{{p{\linewidth}}}]{
                \texttt{\tiny{[P1$\leftrightarrow$P]}}
                \texttt{Post{-}episode Reflexion (not sent to GM):} \\
\texttt{New plan: My previous attempt failed because I assigned a high score (6) to Embarrassment, which misinterprets Mike's reaction to Elise's attack. Mike is not ashamed of his bookstore or his reading habits; rather, he is reacting to the unfairness of the personal attack and the invalidation of his concern. The primary emotion is Defensiveness (protecting his stance and character), followed closely by Anger (at the deflection and insult). Embarrassment should be scored very low (0 or 1) as there is no indication of shame, only a defensive counter{-}argument. Indifference remains 0. Next time, I will strictly differentiate between "feeling attacked" (Defensiveness/Anger) and "feeling ashamed" (Embarrassment). I will assign the highest score to Defensiveness, a high score to Anger, and near{-}zero scores to Embarrassment and Indifference to accurately reflect the intensity of a defensive argument.} \\
            }
        }
    }
    & \\ \\

\end{supertabular}
}

\end{document}
